\section{Os círculos da Clean Architecture em pacotes Quarkus}
\label{sec:circulos-quarkus}

O capítulo~22 do \citetitle{livro:martin2017} desenha quatro círculos
concêntricos e uma regra: as dependências de código-fonte só apontam
para dentro. O círculo de dentro não sabe o nome de nada que está fora
dele, nem de classes, nem de funções, nem de frameworks. A tabela mapeia
cada círculo para pacotes e construções do Quarkus.

\small
\begin{longtable}{@{}p{2.4cm} p{4.4cm} p{5.4cm}@{}}
\toprule
Círculo & Pacote & O que tem e o que pode importar \\
\midrule \endhead
\emph{Entities} & \texttt{\ldots pedidos.dominio} & Records e classes
  com as regras de negócio (\texttt{Pedido}, \texttt{ItemPedido},
  \texttt{Dinheiro}). Só importa \texttt{java.*}. Nada de Jakarta, nada
  de Quarkus \\
\emph{Use Cases} & \texttt{\ldots pedidos.aplicacao} e
  \texttt{\ldots aplicacao.porta} & Casos de uso (\texttt{FecharPedido})
  e as \emph{portas}: interfaces que o caso de uso precisa
  (\texttt{Pedidos}, \texttt{CobrancaGateway}). Importa só
  \texttt{dominio} \\
\emph{Interface Adapters} & \texttt{\ldots adaptador.\allowbreak entrada.\allowbreak rest},
  \texttt{\ldots adaptador.\allowbreak saida.\allowbreak persistencia},
  \texttt{\ldots adaptador.\allowbreak saida.\allowbreak cobranca} & Resources REST (\texttt{@Path}),
  repositórios Panache e entidades JPA, clientes REST. Implementam as
  portas e convertem entre o modelo externo e o domínio \\
\emph{Frameworks \& Drivers} & \texttt{\ldots config} e o próprio
  Quarkus & \texttt{application.properties}, \texttt{@ConfigMapping},
  producers CDI que montam os casos de uso, extensões, banco, Vert.x \\
\bottomrule
\end{longtable}
\normalsize

\paragraph{Onde ficam as extensões.}
Extensões Quarkus são dependências de build, não pacotes seus: entram no
\texttt{pom.xml} e aparecem no código só nos adaptadores e em
\texttt{config}. Uma forma de conferir é olhar os \texttt{import} de
\texttt{dominio} e \texttt{aplicacao}: se aparecer \texttt{io.quarkus},
\texttt{jakarta.persistence} ou \texttt{jakarta.ws.rs}, a regra foi
quebrada. Os adaptadores, por sua vez, importam exatamente esses
pacotes: um resource com \texttt{@Path} ou uma entidade com
\texttt{@Entity} junta, na mesma classe, os dois círculos de fora. O
\citetitle{livro:martin2017} descreve o círculo externo como, em geral,
pouco código além da \emph{cola} que fala com o círculo seguinte, e as
anotações são essa cola. Aceitamos a fusão porque a seta continua
apontando para dentro: o que a regra precisa proteger são os dois
círculos internos, e eles não importam nada do framework.

\paragraph{E a injeção de dependência?}
A pergunta honesta é se \texttt{@ApplicationScoped} e \texttt{@Inject}
podem aparecer em \texttt{aplicacao}. São anotações da especificação
Jakarta (\texttt{jakarta.enterprise.context}, \texttt{jakarta.inject}),
não do Quarkus, e não mudam o comportamento da classe num teste JUnit
puro. Há duas escolhas defensáveis, e a seção seguinte mostra as duas:
\begin{description}
\item[Caso de uso anotado] Mais simples; aceita uma dependência fraca de
  uma especificação estável.
\item[Caso de uso sem anotação, criado por um producer em \texttt{config}]
  Mantém \texttt{aplicacao} com zero dependências externas, ao custo de
  um arquivo a mais. Lembre da seção~\ref{sec:jandex}: classe sem
  \emph{bean defining annotation} não é descoberta, então o producer é
  obrigatório.
\end{description}

O que não é defensável é deixar o domínio depender do framework para
existir: uma entidade de domínio que só funciona dentro de uma transação
JPA, ou uma regra de negócio que lê \texttt{ConfigProvider} diretamente.

\paragraph{Configuração como detalhe.}
O guia \citetitle{doc:quarkus-config-mappings} mostra o
\texttt{@ConfigMapping}: uma interface pública cujos métodos viram
propriedades a partir de um prefixo. Ela fica em \texttt{config}; o caso
de uso recebe só os valores de que precisa (um \texttt{Duration}, um
limite), pelo construtor, sem saber de onde vieram.
